\acresetall
\begin{abstract}
  GPUs are expensive, yet latency-critical \ac{LLM} inference leaves part of them idle.
  \emph{Colocation} reclaims this capacity by running best-effort work, such as fine-tuning, alongside serving.
  Many leading \acp{LLM} now use a \ac{MoE} architecture, which adds a new source of idle time.
  Under expert parallelism, \ac{MoE} experts are spread across GPUs, and because the router sends uneven numbers of tokens to each GPU, lightly loaded GPUs wait for the most loaded one. % at every layer.
  Even after applying load balancing solutions, such \emph{bubbles} still take 8--13\% of prefill time.
  We present \sys, the first colocation system that leverages expert-parallel bubbles for fine-tuning.
  Unlike existing colocation systems, which make one scheduling decision for all GPUs per forward pass, \sys predicts each GPU’s bubble at every layer from the router’s token assignment, and fills it with short chunks of fine-tuning work that reuse the weights already loaded for serving.
  %Unlike existing colocation systems, which make one scheduling decision for all GPUs per forward pass of the entire model, \sys predicts each GPU's bubble duration from the router's decision for each layer and schedules within bubbles short chunks of fine-tuning work of the served weights.
  We also introduce compiler passes that automatically derive these chunks from the served \ac{LLM} as-is, without porting each model architecture to a separate training engine.
  Our evaluation on three \ac{MoE} models, served on four A100 GPUs, shows that fine-tuning throughput with \sys is 1.9--6.2$\times$ faster than with \textsc{LLMStation}, a state-of-the-art colocation system.
  \sys keeps p90 inference latency within 10\% of serving alone in 50 of 63 workloads, compared to 44 for \textsc{LLMStation}. On bursty production traces, \sys trains 3.3$\times$ faster than \textsc{LLMStation} and 14$\times$ faster than a stock trainer colocated under CUDA MPS, with p90 inference latency staying within 7\% of serving alone.

\end{abstract}
